\section{幂等元}\label{sec:4}
\subsection{背景}\label{sec:4a}
先回顾范畴论中的一些初等概念。设 $A$ 是 $\K$ 上的线性范畴，并设某个对象 $Y$ 有幂等自同态 $p\in\Hom_A(Y,Y)$。
\begin{definition*}
对象 $Z$ 连同包含态射与收缩态射 $k,r$，满足
\begin{equation}\label{eq:4.1}
Z\mathrel{\mathop{\rightleftarrows}\limits^{k}_{r}}Y,\qquad rk=e_Z,\quad kr=p,
\end{equation}
称为 $p$ 的\emph{像}。
\end{definition*}
称其为“这个”像是有理由的：给定两个三元组 $(Z,k,r)$、$(Z',k',r')$，两者之间有典范同构 $Z\cong Z'$。
\begin{definition*}
若 $A$ 包含所有幂等自同态的像，就称 $A$ \emph{分裂闭}。
\end{definition*}
特别地，若 $A$ 分裂闭且为加性范畴，那么任意幂等自同态都给出直和分解 $Y=Z\oplus Z^{\perp}$，其中两个直和项分别是 $p$ 与互补投影 $p^{\perp}=e_Y-p$ 的像。
\begin{definition*}
范畴 $A$ 的\emph{分裂闭包}是二元组 $(B,F)$，其中 $B$ 是分裂闭范畴，$F:A\to B$ 是全忠实嵌入，并且 $B$ 的任意对象都是 $A$ 中某个幂等自同态的像。
\end{definition*}
对任意 $A$，都有一个平凡的分裂闭包构造，称为 \emph{Karoubi 完备化} $PA$。\znote{55}{Karoubi 完备化在幂等完备化理论中是标准构造：$PA$ 的对象可等价地写成满足 $p^2=p$ 的二元组 $(Y,p)$，两对象间的态射 $\phi$ 须满足 $\phi=\phi p_0=p_1\phi$，这正是本节公式的另一种写法。嵌入 $A\to PA$ 全忠实，且 $PA$ 的每个对象都是 $A$ 中某个对象的收缩（直和项）。在一切幂等自同态都可分裂的范畴中，$PA$ 是接收 $A$ 的泛对象，故任意分裂闭包都由此得到唯一的等价扩张。[来源：nLab "Karoubi envelope"，https:/\allowbreak /\allowbreak ncatlab.org/\allowbreak nlab/\allowbreak show/\allowbreak Karoubi+\allowbreak envelope（2026-11-15 访问）。]}$PA$ 的对象是二元组 $(Y,p)$，态射空间则是相应的直和项
\[
\Hom_{PA}((Y_0,p_0),(Y_1,p_1))=p_1\Hom_A(Y_0,Y_1)p_0.
\]
嵌入 $A\to PA$ 为 $Y\mapsto(Y,e_Y)$。若 $(B,F)$ 是 $A$ 的任意分裂闭包，则 $F$ 典范地延伸为等价 $PA\to B$。这就证明，任意两个分裂闭包都是等价的，并且这一等价与 $A$
\SourcePage{64}
的有关嵌入相容。当然，在实际工作中，分裂闭包的一般概念多少有些多余，因为通常直接使用 $PA$。\znote{56}{“$A$ 上的三角结构诱导 $PA$ 上的三角结构”是一条经典的稳定性结论：三角范畴的幂等完备化仍是三角范畴，完备化函子 $A\to PA$ 是三角函子，正合三角在完备化后仍正合。证明的关键与本节引理 4.1 相同：把态射沿幂等投影作直和项分解，再验证补全三角的第三个对象仍落在完备化之内。也正因如此，后文才可以在 $PA$ 或 $\Pi\A$ 中自由地取直和项。[来源：nLab "Karoubi envelope"，https:/\allowbreak /\allowbreak ncatlab.org/\allowbreak nlab/\allowbreak show/\allowbreak Karoubi+\allowbreak envelope（2026-11-15 访问）；并见 Seidel, Fukaya categories and Picard–Lefschetz theory, §4a。]}

现在转向与我们的目的尤其相关的三角范畴。可以证明，$A$ 上的三角结构诱导 $PA$ 上的三角结构\bcite{18}。尽管我们不需要这一结果，下面一个相关但更简单的观察将在后文用到：
\begin{lemma}\label{st:4.1}
设 $B$ 是分裂闭三角范畴，$A\subset B$ 是三角全子范畴（即非空，并且对平移与正合三角封闭）。令 $\widetilde A\subset B$ 为如下全子范畴：其对象是所有同构于 $A$ 中某个对象的直和项的对象，故 $\widetilde A\cong PA$。那么 $\widetilde A$ 仍是 $B$ 的三角全子范畴。
\end{lemma}
\begin{proof}
取 $\widetilde A$ 中的态射 $\tilde c:\widetilde Y_0\to\widetilde Y_1$。由定义，每个 $\widetilde Y_j$ 都同构于 $A$ 中某个对象 $Y_j$ 的直和项。将相应的包含态射、收缩态射分别记为 $k_j,r_j$。将 $\tilde c$ 与 $c=k_1\tilde c r_0$ 各自补成正合三角，第三个对象分别记为 $\widetilde Y_2,Y_2$。存在态射 $r_2,k_2$，补全下列交换图：
\SourceDiagram{diagram-4-1}{0.57323015}
其中 $T$ 为平移函子。由于 $j=0,1$ 时 $r_jk_j=\id_{\widetilde Y_j}$，所以 $r_2k_2$ 是 $\widetilde Y_2$ 的自同构（这是三角结构公理的标准推论；若 $A=H^0(\A)$ 来自一个三角化的 $\Ai$ 范畴，我们已将其陈述为\stref{推论}{3.16}）。因此，$\widetilde Y_2$ 是 $Y_2$ 的幂等自同态 $p=k_2(r_2k_2)^{-1}r_2$ 的像，故按定义属于 $\widetilde A$。
\end{proof}

\subsection{幂等元与\texorpdfstring{$\Ai$}{A∞} 模}\label{sec:4b}
以显然的方式，把基域 $\K$ 视为只有一个对象 $*$ 的 $\Ai$ 范畴。当然，它的上同调范畴仍是 $\K$，只是此时将其视为普通范畴。我们在记号上不区分这两者；这应当只会引起极少的混淆。

设 $\A$ 是 $\Ai$ 范畴，$A=H(\A)$ 是其上同调范畴，并令 $\Q=\modcat(\A)$ 为 $\Ai$ 模范畴。
\begin{definition*}
对象 $Y$ 的\emph{同伦幂等元}定义为非含幺 $\Ai$ 函子 $\wp:\K\to\A$，满足 $\wp(*)=Y$。\znote{56}{把基域 $\K$ 视为单对象 $\Ai$ 范畴时，非含幺 $\Ai$ 函子 $\wp:\K\to\A$ 只由作用在唯一对象上的一列元素决定，其相容性方程恰为 (4.2)。在 $\Ai$ 代数的语言中，结构映射 $m_n$ 的次数为 $n-2$（同调次数约定），而 $\wp^d$ 的次数为 $1-d$，两者相差一个平移，这正是函子性方程能逐阶写出的原因。当 $\wp^d=0$（$d\geq2$）时，(4.2) 退化为 $\wp^1\circ\wp^1=\wp^1$，即普通幂等元。[来源：nLab "A-infinity-category"，https:/\allowbreak /\allowbreak ncatlab.org/\allowbreak nlab/\allowbreak show/\allowbreak A-infinity-category（2026-11-15 访问）。]}具体地，它由元素 $\wp^d\in\hom^{1-d}_{\A}(Y,Y)$（$d\geq1$）给出，满足一系列方程
\begin{equation}\label{eq:4.2}
\sum_r\sum_{s_1,\ldots,s_r}\mu^r_{\A}(\wp^{s_1},\ldots,\wp^{s_r})=
\begin{cases}\wp^{d-1},&d\text{ 为偶数},\\0,&d\text{ 为奇数},\end{cases}
\end{equation}
\SourcePage{65}
其中对分拆 $s_1+\cdots+s_r=d$ 求和。
\end{definition*}
\begin{lemma}\label{st:4.2}
设 $p\in\Hom_A(Y,Y)$ 是上同调范畴中的幂等自同态。那么存在同伦幂等元 $\wp$，使 $[\wp^1]=p$。
\end{lemma}
\begin{proof}
$p$ 定义了非含幺函子 $\K\to A$。由\stref{引理}{1.9}，将它提升为 $\Ai$ 函子的障碍位于 $\HH^2(\K,A)^{2-r}$ 中，其中 $r>2$。基础 Hochschild 复形为
\[
\begin{gathered}
CC^{r+s}(\K,A)^s=\Hom_A^s(Y,Y),\\
\partial^{r,s}(T)=(-1)^s
\begin{cases}pT-Tp,&r\text{ 为偶数},\\pT-Tp^{\perp},&r\text{ 为奇数}.\end{cases}
\end{gathered}
\]
将 $\Hom_A(Y,Y)$ 按照与 $p$ 作左复合、右复合的像与核进行分解，就可直接看出，唯一非零的上同调群为
\[
\HH^s(\K,A)^s=p\Hom_A^s(Y,Y)p\oplus p^{\perp}\Hom_A^s(Y,Y)p^{\perp}.
\]
\end{proof}
\begin{note}\label{st:4.3}
除了直接计算 Hochschild 复形，也可以使用一些基本的同调代数。先向 $\K$ 添加一个单位元，使所讨论的函子变为含幺的；这样，\stref{笔记}{2.7}中的相应双模范畴就是半单的，因此所有高阶 $\operatorname{Ext}$ 群都为零。
\end{note}
设 $\wp$ 是同伦幂等元。与之相应，我们定义一个 $\Ai$ 模 $\mathcal Z$，称为 $\wp$ 的\emph{抽象像}：
\begin{definition*}

\[
\mathcal Z(X)=\hom_{\A}(X,Y)[q]
\]
是以 $\hom_{\A}(X,Y)$ 为系数、以分次变量 $q$（$|q|=-1$）为变量的形式多项式空间。用 $\delta_q$ 表示正规化的形式微分运算：它将 $q^0$ 映为零，将其他幂 $q^k$ 映为 $q^{k-1}$。用这一记号，微分为
\begin{equation}\label{eq:4.3}
\begin{aligned}
\mu^1_{\mathcal Z}(b(q))={}&\sum_r\sum_{s_2,\ldots,s_r}\delta_q^{s_2+\cdots+s_r}\mu^r_{\A}(\wp^{s_r},\ldots,\wp^{s_2},b(q))\\
&+(-1)^{|b|}\frac{b(q)-b(-q)}{2q},
\end{aligned}
\end{equation}
其中最后一项是先将幂级数反对称化，再施加 $\delta_q$；它实际上不涉及乘以 $\frac12$，因此在 $\K$ 的特征为二时也有意义。对于 $d\geq2$，高阶结构映射为
\[
\begin{aligned}
&\mu^d_{\mathcal Z}(b(q),a_{d-1},\ldots,a_1)\\
&\quad=\sum_r\sum_{s_{d+1},\ldots,s_r}\delta_q^{s_{d+1}+\cdots+s_r}
\mu^r_{\A}(\wp^{s_r},\ldots,\wp^{s_{d+1}},b(q),a_{d-1},\ldots,a_1).
\end{aligned}
\]
\end{definition*}
\SourcePage{66}
\begin{lemma}\label{st:4.4}
基础上同调模 $H(\mathcal Z)$ 可以如下描述：$H(\mathcal Z)(X)\cong p\Hom_A(X,Y)$，其中 $p=[\wp^1]\in\Hom_A(Y,Y)$；模结构 $H(\mathcal Z)(X_1)\otimes\Hom_A(X_0,X_1)\to H(\mathcal Z)(X_0)$ 来自通常的复合
\[
p\Hom_A(X_1,Y)\otimes\Hom_A(X_0,X_1)\longrightarrow p\Hom_A(X_0,Y).
\]
\end{lemma}
\begin{proof}
$\mathcal Z(X)$ 有一个自然的递增滤过，即 $F_r=\hom_{\A}(X,Y)q^0\oplus\cdots\oplus\hom_{\A}(X,Y)q^{r-1}$。我们通常使用递降滤过；为统一起见，将 $r$ 换成 $-r$，就得到一个递降滤过。相应谱序列满足
\[
E_1^{-r,s}=H^{s-r}(F_{r+1}/F_r)=
\begin{cases}\Hom_A^s(X,Y),&r\geq0,\\0,&\text{其余情形}.\end{cases}
\]
微分由\eref{4.3}中恰含一个 $\delta_q$ 或 $1/q$ 的项诱导，因此为
\[
\partial^{-r,s}(b)=(-1)^{r+s}
\begin{cases}p^{\perp}b,&r>0\text{ 且为奇数},\\pb,&r>0\text{ 且为偶数},\\0,&\text{其余情形}.\end{cases}
\]
所以谱序列在 $E_2$ 页退化，$H(\mathcal Z)(X)$ 同构于唯一非零的列 $E_2^{0\bullet}=p\Hom_A(X,Y)$。事实上，谱序列的边缘同态给出这两个群之间的典范同构；利用这一同构，所需的模结构相容性就是显然的。
\end{proof}
\begin{lemma}\label{st:4.5}
对任意 $\Ai$ 模 $\M$，有同构
\[
\Hom_{H(\Q)}(\mathcal Z,\M)\cong H(\M)(Y)p,
\]
并且它对于以模同态作左复合而言是自然的。
\end{lemma}
\begin{proof}
考虑\stref{笔记}{2.15}中的 Abel 范畴 $\operatorname{Mod}^0(A)$。\stref{引理}{4.4}说明，作为该范畴的对象，$H(\mathcal Z)$ 是 $\Hom_A(-,Y)$ 的直和项；后者为投射对象，故前者也是。结合\eref{2.14}，这意味着 $\hom_{\Q}(\mathcal Z,\M)$ 上长度滤过所诱导的谱序列在 $E_2$ 页退化，而且映射
\[
\Hom_{H(\Q)}(\mathcal Z,\M)\longrightarrow\Hom_{\operatorname{Mod}^0(A)}(H(\mathcal Z),H(\M)[\bullet])
\]
是同构。由经典 Yoneda 引理及\stref{引理}{4.4}，后一个分次群自然同构于 $H(\M)(Y)p$。
\end{proof}
考虑将 Yoneda 嵌入应用于 $p$ 所得的 $H^0(\Q)$ 中的幂等自同态。由\stref{引理}{4.5}直接可知，$\mathcal Z$ 是这个自同态在\eref{4.1}意义下的像。特别地，在 $H^0(\Q)$ 中的同构意义下，抽象像只依赖于 $p$，而不依赖于 $\wp$ 的选择。
\SourcePage{67}
\begin{note}\label{st:4.6}
$\mathcal Z$ 的定义受到\bcite{24}中一个观察的启发：在适当的复形范畴中，利用无限分解的技巧可以得到直和项的存在性（感谢 François Loeser 向我指出这一点）。
\end{note}

\subsection{分裂闭的\texorpdfstring{$\Ai$}{A∞} 范畴}\label{sec:4c}
\begin{definition*}
若任意同伦幂等元的抽象像都由 $\A$ 的某个对象拟表示，就称 $\A$ \emph{分裂闭}。$\Ai$ 范畴 $\A$ 的\emph{分裂闭包}是二元组 $(\B,\F)$，其中 $\B$ 为分裂闭 $\Ai$ 范畴，$\F:\A\to\B$ 为上同调全忠实函子，且在 $H^0(\B)$ 中，每个对象都是 $H^0(\A)$ 中某个对象的幂等自同态的像。
\end{definition*}
表述这些概念时，我们沿用了\sr{sec:3}{第 3 节}的一般模式，但其实可以大幅简化它们的含义。首先，$\A$ 分裂闭，当且仅当 $H^0(\A)$ 具有这一性质（按最初在上同调层次上的含义）。事实上，设 $H^0(\A)$ 分裂闭，并取同伦幂等元 $\wp$，满足 $[\wp^1]=p$。由假设，$p$ 在 $H^0(\A)$ 中有像 $Z$。利用\stref{引理}{4.5}，很容易证明 $Z$ 的 Yoneda 模在 $H^0(\Q)$ 中同构于 $\wp$ 的抽象像，因此给出了一个拟表示对象。逆向蕴含类似，但还需使用\stref{引理}{4.2}。直接推论是：$(\B,\F)$ 为 $\A$ 的分裂闭包，当且仅当 $(H^0(\B),H^0(\F))$ 为 $H^0(\A)$ 的分裂闭包。
\begin{lemma}\label{st:4.7}
每个 $\Ai$ 范畴都有分裂闭包。此外，若 $(\B,\F)$ 与 $(\widetilde{\B},\widetilde{\F})$ 是 $\A$ 的两个分裂闭包，则存在拟等价 $\G:\B\to\widetilde{\B}$，使 $\G\circ\F$ 在 $H^0(\operatorname{fun}(\A,\widetilde{\B}))$ 中与 $\widetilde{\F}$ 同构。
\end{lemma}
\begin{proof}
用 $\Pi\A$ 表示 $\Q=\modcat(\A)$ 的如下全 $\Ai$ 子范畴：其对象是在 $H^0(\Q)$ 中同构于某个同伦幂等元 $\wp$ 的抽象像 $\mathcal Z$ 的所有对象。我们断言，$\Pi\A$ 包含所有与 $\A$ 的对象 $Y$ 相应的 Yoneda $\Ai$ 模 $\Y$。为说明这一点，取同伦幂等元，使 $[\wp^1]$ 为 $Y$ 的恒等自同态（由\stref{引理}{4.2}，它存在）；由\stref{引理}{4.5}，$\wp$ 的抽象像同构于 $Y$ 的 Yoneda 模。于是二元组 $(\Pi\A,\mathcal I_{\A}:\A\to\Pi\A\subset\Q)$ 就是 $\A$ 的分裂闭包。事实上，如前所述，只需考察上同调范畴；在那里，同样结合\stref{引理}{4.2}与\stref{引理}{4.5}就得到所需结论。

对于第二个结论，设 $(\B,\F)$ 是 $\A$ 的任意分裂闭包。利用 $\B$ 的 Yoneda 嵌入与 $\Ai$ 模的拉回，定义 $\Ai$ 函子
\[
\G:\B\xrightarrow{\mathcal I_{\B}}\modcat(\B)\xrightarrow{\F^*}\Q.
\]
由于 $\F$ 上同调全忠实，$\G\circ\F$ 同构于 $\mathcal I_{\A}$，见\eref{2.13}。这意味着，将 $\G$ 限制在形如 $\F(Y)$ 的对象上时，它上同调全忠实；但由定义，$H^0(\B)$ 的每个对象都是某个 $\F(Y)$ 的幂等
\SourcePage{68}
自同态的像，因此 $\G$ 在整个 $\B$ 上都上同调全忠实。类似论证表明，$\G$ 的像包含于 $\Pi\A$，而且 $\G$ 实际上是拟等价 $\B\to\Pi\A$。当然，既然已有任意分裂闭包与 $\Pi\A$ 之间的拟等价，就可以利用\stref{定理}{2.9}构造任意两个这样的闭包之间的拟等价。
\end{proof}
\begin{lemma}\label{st:4.8}
三角化的 $\Ai$ 范畴的分裂闭包仍然三角化。
\end{lemma}
\begin{proof}
设 $\A$ 为三角化的 $\Ai$ 范畴。由\stref{例题}{3.13}，$\Q=\modcat(\A)$ 三角化；经 Yoneda 嵌入，$H^0(\A)$ 被识别为 $H^0(\Q)$ 的一个经典意义下的三角子范畴。此外，$H^0(\Pi\A)$ 是由同构于 $H^0(\A)$ 中某个对象的直和项的对象组成的子范畴，因此由\stref{引理}{4.1}，它仍是 $H^0(\Q)$ 的三角子范畴。换言之，若 $H(\Q)$ 中有正合三角，其中两个对象属于 $\Pi\A$，那么第三个也属于 $\Pi\A$。由\stref{推论}{3.10}，这意味着 $\Pi\A$ 本身是三角化的 $\Ai$ 范畴。
\end{proof}
\begin{definition*}
设 $\A$ 是任意 $\Ai$ 范畴；先取它的三角化包络，再取后者的分裂闭包，将结果记为 $\B$。例如，两次都采用标准选择，就得到 $\B=\Pi(\Tw\A)$。称 $H^0(\B)$ 为\emph{分裂闭导出范畴}，记为 $D^{\pi}(\A)$。
\end{definition*}
与\secref{3j}的讨论相对应，还应提到一个相关概念。设 $\B$ 是分裂闭、三角化的 $\Ai$ 范畴，$\A\subset\B$ 为非空的全 $\Ai$ 子范畴。考虑包含 $\A$、对 $H^0(\B)$ 中的同构封闭、且本身三角化并分裂闭的最小全子范畴 $\widetilde{\B}\subset\B$。
\begin{definition*}
称它为 $\B$ 中由 $\A$ \emph{分裂生成}的子范畴。\znote{60}{分裂生成在 §5 中反复出现：推论 5.7 断言，若某族对象使所有 $X$ 的逐次扭转都为零，则该族既生成 $\A$、又使 $\A$ 分裂闭；推论 5.8 则处理 $T_{Y_1}\cdots T_{Y_m}X\cong S^{\sigma}X$（$\sigma\ne0$）的情形，结论是 $Y_k$ 分裂生成 $\A$。直观上，分裂闭包把“由映射锥生成”与“取直和项”两件事一并完成，这正是 Fukaya 范畴中由球面等对象生成整个范畴时出现的现象。[来源：并见 Seidel, Fukaya categories and Picard–Lefschetz theory, §5e（推论 5.7 与 5.8）。]}
\end{definition*}
具体地，从 $\Ob\A$ 出发，取映射锥和平移，任意多次地重复这一过程，最后取先前构造的对象的所有幂等自同态的像，就得到 $\Ob\widetilde{\B}$。若 $\widetilde{\B}=\B$，就称 $\A$ \emph{分裂生成} $\B$，或称 $\A$ 的对象是 $\B$ 的\emph{分裂生成元}。例如，$\Pi(\Tw\A)$ 由 $\A$ 的对象分裂生成。由\stref{引理}{3.20}与\stref{引理}{4.7}中的唯一性结论，立即得到：
\begin{corollary}\label{st:4.9}
设 $\B$ 是分裂闭、三角化的 $\Ai$ 范畴，$\A\subset\B$ 为分裂生成 $\B$ 的全子范畴。那么存在拟等价 $\Pi(\Tw\A)\to\B$，它诱导三角范畴的等价 $H^0(\B)\cong D^{\pi}(\A)$。
\end{corollary}
\begin{note}\label{st:4.10}
分裂闭包 $\Pi\A$ 有一点在美学上不够理想：它并不把 $\A$ 本身包含为全子范畴（当然，它包含一个与 $\A$ 拟同构的子范畴）。现在概述另一种构造，它是无界扭曲复形的一种变体。这一构造更
\SourcePage{69}
具体，也更接近最初的想法；但至少这里给出的版本需要严格含幺。

设 $\A$ 是严格含幺的 $\Ai$ 范畴，并假设它有零对象 $O$，使 $\hom_{\A}(O,X)$ 与 $\hom_{\A}(X,O)$ 都为零群。如下定义另一个 $\Ai$ 范畴 $\Sigma^-\A$。其对象是形式无限序列 $X=(X^0,X^{-1},X^{-2},\ldots)$，其中 $X^i\in\Ob\A$。态射 $a\in\hom^k_{\Sigma^-\A}(X_0,X_1)$ 是无限矩阵
\[
a=(a^{ji}),\qquad a^{ji}\in\hom^{k+j-i}_{\A}(X_0^i,X_1^j),
\]
要求它“本质上下三角”：存在某个 $N$（依赖于 $a$），使 $j<i-N$ 时 $a^{ji}=0$。复合映射通过矩阵乘法定义，与加性扩张的情形相同。注意，每个 $X^i$ 部分都被视为平移了 $i$，这会产生一些符号。事实上，为简化下面的公式，我们采用\stref{笔记}{3.26}中修改后的符号约定。

考虑二元组 $(X,\delta_X)$，其中 $X\in\Ob\Sigma^-\A$，$\delta_X\in\hom^1_{\Sigma^-\A}(X,X)$ 满足\eref{3.19}，并且严格下三角，即 $j\leq i$ 时 $\delta_X^{ji}=0$。定义 $\Ai$ 范畴 $\Tw^-\A$，其对象为这样的二元组，态射的结构映射则像\eref{3.20}那样用微分进行变形。可以把 $\A$ 与 $\Tw^-\A$ 中由序列 $(X^0,X^{-1}=O,X^{-2}=O,\ldots)$ 组成的全 $\Ai$ 子范畴等同。

给定 $Y\in\Ob\A$ 的同伦幂等元 $\wp$，如下定义 $\Tw^-\A$ 的对象 $(Z,\delta_Z)$：
\begin{equation}\label{eq:4.4}
\begin{gathered}
Z^i=Y\quad\text{对所有 }i,\\
\delta_Z^{ji}=\begin{cases}
\wp^1,&i\text{ 为偶数且 }j=i+1\leq0,\\
\wp^1-e_Y,&i\text{ 为奇数且 }j=i+1\leq0,\\
\wp^{j-i},&i+1<j\leq0,\\
0,&\text{其余情形}.
\end{cases}
\end{gathered}
\end{equation}
广义 Maurer–Cartan 方程等价于\eref{4.2}。定义
\[
\begin{aligned}
r&\in\hom^0_{\Tw^-\A}(Y,Z),&r^{00}&=e_Y,\\
k&\in\hom^0_{\Tw^-\A}(Z,Y),&k^{0i}&=\wp^{1-i},\\
h&\in\hom^{-1}_{\Tw^-\A}(Z,Z),&h^{i-1,i}&=(-1)^i e_Y,
\end{aligned}
\]
除列出的系数外，其余系数都为零。于是
\[
\begin{aligned}
\mu^1_{\Tw^-\A}(r)&=0,&\mu^1_{\Tw^-\A}(k)&=0,\\
\mu^2_{\Tw^-\A}(k,r)&=\wp^1,&\mu^2_{\Tw^-\A}(r,k)-\mu^1_{\Tw^-\A}(h)&=e_Z.
\end{aligned}
\]
\SourcePage{70}
这意味着，在上同调范畴中，$Z$ 是 $Y$ 的幂等自同态 $[\wp^1]$ 的像。因此，若取 $\Tw^-\A$ 中由 $\A$ 的所有对象以及所有\eref{4.4}类型的对象组成的全 $\Ai$ 子范畴，就得到 $\A$ 的一个分裂闭包。为说明它与前一种方法的联系，注意与 $Z$ 相应的 Yoneda $\A$ 模恰好就是抽象像 $\mathcal Z$。
\end{note}
